\section{Absolute Zero: dejar que el modelo proponga su propia siguiente pregunta}

\subsection{La propia tarea de entrenamiento se vuelve el cuello de botella de datos}

\videofigure{figures/task-data-bottleneck.jpg}{El reasoning training actual depende de que expertos humanos escriban el prompt, la respuesta o el reasoning trace; entre más fuerte es la capacidad, más difícil resulta encontrar suficientes tareas de frontera.}{00:23:46--00:24:38}

El verifier del RLVR tradicional puede calificar automáticamente, pero el conjunto de problemas sigue siendo provisto de forma estática por humanos. Cuando el modelo progresa, las preguntas simples pierden gradiente y las demasiado difíciles se fallan por completo; el entrenamiento necesita mantenerse constantemente en la current frontier de tareas.

\videofigure{figures/proposer-solver.jpg}{Absolute Zero usa el mismo modelo para desempeñar simultáneamente los papeles de proposer y solver, generando tareas por sí mismo y resolviéndolas él mismo dentro de un entorno de código ejecutable.}{00:24:38--00:25:24}

\videofigure{figures/absolute-zero-loop.jpg}{El proposer genera tareas según el learnability reward, el solver resuelve los problemas según el accuracy reward, y ambas señales actualizan conjuntamente el modelo.}{00:25:24--00:26:18}

\begin{importantbox}{Lo que Absolute Zero elimina es el task data externo, no el entorno}
El sistema sigue necesitando un ejecutor de código que proporcione retroalimentación verificable. Lo que se llama zero data se refiere a no depender de un conjunto de problemas redactado de antemano por humanos, no a aprender en un vacío sin entorno ni reward.
\end{importantbox}

\subsection{Deduction, Abduction e Induction}

\videofigure{figures/proposer-modes.jpg}{El proposer amplía su currículo autogenerado mediante tres modos de tareas de código: deduction, abduction e induction.}{00:26:18--00:27:10}

Los tres tipos de tareas pueden entenderse así:

\begin{itemize}
    \item \textbf{Deduction}: dados un programa y una entrada, predecir la salida de la ejecución;
    \item \textbf{Abduction}: dados un programa y una salida, deducir en reversa una entrada que satisfaga la condición;
    \item \textbf{Induction}: dados varios pares de entrada y salida junto con una descripción en lenguaje natural, inferir el programa o la regla.
\end{itemize}

Estas tareas pueden ser evaluadas por el ejecutor, por lo que son adecuadas para generarse, resolverse y verificarse de manera automática; además, las tres direcciones exigen distintos tipos de reasoning, lo que ayuda a evitar que el currículo sea demasiado uniforme.

\subsection{Learnability reward: ni demasiado fácil, ni completamente imposible}

Sea $\bar p$ la tasa de éxito empírica del solver en una misma tarea. La versión simplificada del proposer reward descrita en el curso es:

$$
r_{\mathrm{learn}}=
\begin{cases}
0, & \bar p=0,\\
1-\bar p, & \bar p>0.
\end{cases}
$$

\begin{itemize}
    \item $\bar p=1$: la tarea ya se domina, la learnability es baja;
    \item $\bar p=0$: la tarea probablemente es inválida o excede con mucho la capacidad actual, tampoco se otorga recompensa;
    \item $0<\bar p<1$: la tarea se encuentra en el límite de la capacidad, con muestras exitosas y a la vez margen de mejora.
\end{itemize}

\begin{knowledgebox}{Esto es curriculum learning automatizado}
A medida que el solver se vuelve más fuerte, las tareas que antes eran difíciles se vuelven simples, y el proposer debe aumentar continuamente la complejidad. El currículo ya no está definido por un banco de preguntas fijo, sino que se genera dinámicamente a partir del límite actual de fallas del modelo.
\end{knowledgebox}

\subsection{Task validation y buffer management}

\videofigure{figures/task-validation.jpg}{Antes de entrar al entrenamiento, las tareas deben pasar comprobaciones de ejecutabilidad, seguridad y determinismo, y se escriben en un task buffer con estadísticas de éxito.}{00:28:10--00:29:28}

La validación incluye verificar si el programa puede ejecutarse, si desencadena operaciones peligrosas, y si la misma entrada produce de manera determinista la misma salida. El buffer guarda el programa, la entrada, la salida, la tasa de éxito y tareas de referencia, de modo que el proposer pueda tomar muestras del currículo histórico y generar nuevas tareas distintas pero aprendibles.

\begin{warningbox}{El task proposer también puede hacer reward hack}
Puede generar preguntas que exploten vulnerabilidades del intérprete, preguntas disfrazadas de difíciles pero en realidad carentes de sentido, o repetir tareas antiguas ligeramente reescritas. Se necesitan sandbox, deduplicación, restricciones de complejidad y validación independiente.
\end{warningbox}

\subsection{Resultados y preguntas abiertas}

\videofigure{figures/absolute-zero-takeaways.jpg}{Absolute Zero logra un reasoning de código sólido sin datos de prompt elaborados por humanos, y presenta un emergent curriculum en el que la complejidad y la diversidad de las tareas crecen con el entrenamiento.}{00:29:28--00:33:34}

El curso subraya que el proposer y el solver mantienen una ligera relación adversarial: el proposer busca encontrar problemas que el solver todavía no sabe resolver pero puede aprender, mientras que el solver comprime constantemente esa región. Los resultados muestran que las tareas autogeneradas pueden superar a la línea base entrenada con una gran cantidad de ejemplos elaborados por expertos.

\begin{importantbox}{La verdadera apertura proviene de que la distribución de tareas también puede cambiar}
Repetir RL únicamente sobre un banco de preguntas fijo, en el mejor de los casos se aproxima al límite superior de esa distribución; solo dejar que el modelo cambie el problema mismo abre la posibilidad de seguir ampliando el límite del aprendizaje. Pero la validez de las tareas y su relevancia con la realidad se convierten así en un nuevo problema para el verifier.
\end{importantbox}

\subsection{Resumen de la sección}

Absolute Zero lleva el synthetic data de "generar uno mismo las respuestas" a "generar uno mismo tanto las preguntas como las respuestas". El entorno de ejecución de código proporciona verificación instantánea, lo que permite que el par proposer--solver forme un currículo dinámico; al extenderlo hacia dominios abiertos del mundo real, la fidelidad del entorno y el reward se convierten en las limitaciones centrales.
